\chapter{General discussion} \label{ch: general discussion}

The studies presented in the preceding chapters present exploratory yet systematic analyses of the nature of inter-individual variability in the encoding and decoding of prosodic prominence relations. Two production studies and one perception study investigated individual differences in the mapping of categorical and continuous prosodic cues onto functional aspects of prominence. This chapter discusses the empirical and theoretical contributions of these studies.

\section{Summary of results}

To begin with, this section summarizes the main findings of all three studies. In the first production study presented in Chapter \ref{ch:Production I}, the marking of newness versus accessibility in two subsequent referents was investigated. Results indicated that speakers were rather similar in the type of continuous prosodic parameters used to mark newness, yet differed in their strength of encoding. A cluster analysis revealed a split into two groups, the weak encoders and the strong encoders. In terms of the categorical prosodic parameters, one cluster of speakers was slightly more likely to use accent status as a cue to information status, the other was more likely to use accent type, indicating a weak trade-off effect.

In the second production study presented in Chapter \ref{ch:Production II}, the effects of lexical and referential aspects of information status on prosodic prominence were analyzed separately and the syntactic factors sentence position and grammatical role were additionally included. Individual variability in the encoding of lexical information status was again mostly related to the strength of prosodic marking, while there was more variability in the encoding of the referential level: All speakers used higher periodic energy mass in r-new words, but one group of speakers additionally relied on F0 contour shape and alignment, while the other group used F0 scaling. There were even more diverse strategies in the prosodic marking of grammatical role and sentence position, which mostly related to strength of encoding but also to the type of cue used. In addition, there was a trade-off between the use of accent status and accent type similar to the first study, particularly in the encoding of referential givenness, grammatical role, and sentence position. 

This second production study also considered communicative skill as a potential factor towards explaining individual variability in prosodic form-meaning mapping. As expected, individuals with higher communicative skill produced larger prosodic differences between l-new and l-given referents, especially in terms of accent status, duration, and intensity. However, unexpectedly, the reverse relation was observed for referential information status: Speakers with \textit{lower} communicative skill produced larger differences in terms of accent type and F0 scaling between r-new and r-given words. The results for grammatical role were similar to those for lexical information status, while the prosodic marking of sentence position barely seemed to be influenced by communicative skill.

In the perception study presented in Chapter \ref{ch: Perception}, listeners were asked to rate the perceived prominence in the target words investigated in Production Study II. We found accent status to be the most important cue to prominence followed by duration and accent type, while non-prosodic parameters played only a small role. Individuals did not vary much in their cue rankings, but one group of listeners emerged that was more sensitive to the prominence cues than the other group. However, communicative skill offered little to explain this difference. Only the differentiation of accent types seemed to improve in individuals with higher communicative skill.

\section{Variability in strength of prominence encoding and decoding effects}

A pattern that was repeatedly found throughout all experiments is that a smaller group of speakers or listeners more strongly marks meaning contrasts via prosody or is more sensitive to prominence cues in perception, while a larger group shows the same effects in a more subtle way. This is in line with previous findings, for example, by \citet{roessigDynamicsIntonationCategorical2019}, who observed that a larger group of speakers produced smaller distinctions between focus types in terms of tonal onglide and a smaller group produced larger differences. In general, such findings may help explain why prosodic effects in empirical studies averaged across multiple individuals are often subtle and sometimes inconsistent. Thinking of strength of prosodic encoding as a continuum, where speakers occupy different places between the endpoints ``no encoding" and ``strong encoding", a reasonable relation can be drawn to Lindblom's \citeyearpar{lindblomExplainingPhoneticVariation1990} classic H\&H theory. Specifically, strength of encoding may be related to different weightings of the hyperarticulation and hypoarticulation principles. A higher weighting of the hyperarticulation principle thereby reflects higher listener orientation \citep{turnbullListenerorientedPhoneticReduction2019} and one way this may manifest is in a stronger prosodic marking of meaning-related contrasts. Assuming that communicative skill is a correlate of higher attentiveness towards listener needs, we see a confirmation of this general principle in Production Study II: Higher communicative skill correlates with stronger prosodic differentiation of lexical newness and givenness as well as grammatical object and subject function. 

These considerations would suggest that some speakers are better communicators than others, as has been suggested, for example, by \citet{cangemiListenerSpecificPerceptionSpeakerSpecific2015}. However, data presented in Production Study II indicate that this generalization is too simple. Instead of being generally good or bad communicators, individuals may just differ in how well they encode specific aspects of meaning. This is evident in the clustering results, where strong encoders of lexical information status overlap only marginally with strong encoders of referential givenness. Furthermore, the converse effects of communicative skill on lexical and referential information status marking suggest that different skill sets facilitate the encoding of different meaning contrasts.

Communicative skill could also be interpreted as a superior ability to weigh listener needs (Lindblom's hyperarticulation) against required effort in production (Lindblom's hypoarticulation). Similar to the H\&H theory, the Smooth Signal Redundancy Hypothesis (SSRH) argues that speakers weigh language redundancy against prosodic prominence in order to achieve robust communication \citep{aylettSmoothSignalRedundancy2004}. Interpreted in this light, the weaker prosodic encoding of referential givenness in speakers with higher communicative skill could be explained by their awareness of the presence of another cue, morphosyntactic marking, which encourages prosodic hypoarticulation. That prosodic markers may be weighed against non-prosodic factors is also in line with a finding by \citet{giovannoneIndividualDifferencesLexical2023} that listeners are able to trade off lexical and acoustic information in speech perception.

An alternative explanation relates directly to the communicative situations established in the production studies: Speakers were asked to produce texts in a listener-oriented manner, so that their presumed fellow participants could then fulfil a comprehension task. This set-up could encourage speakers with higher communicative skill to hyperarticulate every content word in order to help the listener memorize the relevant information, similar to a child-directed speech style \citep[e.g.,][]{cristiaInputLanguagePhonetics2013}. We see traces of such behavior in individual speakers that frequently produce phrase boundaries after target words or L+H* accents across conditions, for example, s31 in Production Study I or s08 in Production Study II. The failure to prosodically differentiate r-new from r-given referents would thus result from both words being produced highly prominently. However, our data indicate that r-new referents are rather produced with reduced prominence in such cases, not providing sufficient evidence for this potential explanation. 

In perception, communicative skill has initially been interpreted as an enhanced ability to interpret pragmatic information, i.e., speaker-intended meaning \citep{nieuwlandIncrementalityPragmaticProcessing2010, bishopIndividualDifferencesTopdown2016}. However, recent research has also postulated a relation between communicative skill and sensitivity to prosodic differences \citep{bishopPhonologyPhoneticsSignalextrinsic2020, orricoInfluenceEmpathyAutisticlike2023}. In the perception study presented in Chapter \ref{ch: Perception}, we did not observe any correlation between communicative skill and reliance on pragmatic information in prominence ratings. In fact, pragmatic cues were among the least important predictors of prominence in all listeners. We did however observe slightly different degrees of sensitivity to differences between pitch accent types, which were modulated by communicative skill. More communicatively skilled listeners were able to make more fine-grained distinctions between accent types. 

Since our prominence rating task did not specifically reference meaning distinctions, it is perhaps unsurprising that we neither observed any strong overall effects of non-prosodic factors nor individual differences in pragmatic reliance \citep[cf.][]{coleListeningSoundListening2014, coleSoundStructureMeaning2019}. In addition, the communicative skill score collected via the AQ questionnaire is somewhat broad and it is unclear to what extent it captures pragmatic skill or prosodic sensitivity. There is a remaining need to tease apart these potentially different aspects of individual communicative skill.

Varying strength of effects observed across individuals in all studies can also be related to the degree of engagement with the task or differences in metalinguistic awareness necessary to perform the tasks. Especially the prominence rating task employed in our perception study explicitly taps into metalinguistic awareness and may thus be subject to such effects. \citet{niebuhrF0basedRhythmEffects2009} argues that more musical individuals are particularly skilled in performing these tasks. However, a musical advantage does not necessarily translate into a linguistic processing benefit \citep{niebuhrF0basedRhythmEffects2009}. More indirect, communicatively meaningful tasks may thus yield additional insights into individual variability in the level of encoding and decoding abilities.

\section{Cue trading and prosodic categories}

Compared to the extent of variability observed with regard to strength of encoding, there seems to be relatively little variability across speakers and listeners in the type of cues that they use. Especially in the perception study, listeners were rather uniform with regard to the importance rankings of prominence cues, with accent status clearly being the most important predictor of prominence ratings for all. This uniformity may be somewhat unexpected considering that trade-off or cue weighting effects, which predominantly relate to the variability regarding the use of different types of (continuous) prosodic parameters, have recently become the primary interest of studies specifically investigating inter-individual differences \citep[e.g.,][]{arvanitiVariabilityOverlapCue2024, severijnenYourVOORnaamNot2024}. 

Nevertheless, we observe some cue trading effects in production between F0 scaling and alignment in Production Study II and accent type and accent status in both Production Study I and II. The latter effect is particularly interesting as it relates to the independency of tune and metrical structure, one of the premises of AM theory \citep{laddIntonationalPhonology2008}. Previous observations regarding cue trading effects have mostly been made in relation to continuous prosodic cues \citep{ipSearchSalienceFocus2021, arvanitiVariabilityOverlapCue2024, severijnenYourVOORnaamNot2024}, although prosodic categories are arguably decisive for the interpretation of intonational meaning \citep{laddIntonationalPhonology2008}. The findings in the present production studies thus add new insights on the nature of individual variability and the possibility of cue trading on a categorical level.

Production Study I also revealed that stronger encoding via continuous phonetic cues does not necessarily correlate with a more consistent use of categorical parameters. Specifically, the speakers that showed stronger effects of information status on continuous measures of F0 were less likely to vary accent status as a function of this meaning contrast. Since accent status was confirmed as the most important perceptual cue to prominence in our perception study \citep[also in line with previous studies by][]{baumannWhatMakesWord2018, wagnerExploitingSpeechgestureLink2019}, this finding further stresses the importance of taking into account categorical prosodic parameters. A lack of correspondence between continuous and categorical parameters could be attributed to the pervasiveness of individual differences in the differentiation of prosodic categories. Even a strong categorical contrast like accentuation is subject to variability. For example, \citet{beckmanArticulatoryEvidenceDifferentiating1994} report that for some but not all speakers nuclear accented syllables are longer than unaccented syllables.

In general, all studies confirm the prevalence of prosodic categories for prominence production and especially perception not just on an overall level but also for the individual. Prosodic categories, such as accent status or pitch accent type, may be especially robust to individual variability as they represent bundles of cues. Listeners with diverse preferences for prosodic cues may still be able to interpret and make use of them as they combine several acoustic dimensions. A similar argument has been advanced by \citet{bruggemanUnifyingSpeakerVariability2017} for TCoG as a cue bundle that more reliably correlates with focus type across speakers of Egyptian Arabic than more isolated F0 measures based on turning points.

\section{Relationality, redundancy, and robustness of communication}

Spoken language has been described as a robust and redundant system \citep{winterSpokenLanguageAchieves2014}, in analogy to a theoretical model in biology \citep{kitanoBiologicalRobustness2004}.\footnote{\citet{winterSpokenLanguageAchieves2014} distinguishes between the terms \textit{redundancy} and \textit{degeneracy}. Redundancy refers to repetitions of the same component, while degeneracy refers to different components performing the same function. Here, we forgo this distinction and use the term redundancy throughout, in line with \citet{jasminTailoredPerceptionIndividuals2020}.} Redundancy refers to the concept of multiple components serving the same function, ensuring the robustness of the system in case of the failure of one component \citep{kitanoBiologicalRobustness2004, winterSpokenLanguageAchieves2014}. Robustness in speech can be understood as optimal transmission of linguistic information resulting in successful communication, especially in the face of noise or variation \citep{winterSpokenLanguageAchieves2014}.

Prosodic prominence specifically can be regarded as highly redundant. Prominence can be expressed by single or multiple cues related to the acoustic dimensions F0, duration, and intensity \citep{kochanskiLoudnessPredictsProminence2005, baumannWhatMakesWord2018, roessigTracingPhoneticSpace2022}, ensuring signal redundancy and thus communicative robustness. Furthermore, the relational nature of prominence contributes to this robustness and allows for some variability in the realization of individual prosodic constituents \citep{terkenFundamentalFrequencyPerceived1991, terkenPerceptionProsodicProminence2000, laddIntonationalPhonology2008}. As long as general prominence relations are preserved, prominence is adequately encoded by the speaker and decodable for the listener. We see evidence for this in Production Study II, where lexically and referentially given target words followed by a nuclear accented verb are variably produced with a prenuclear or no accent, so that a weak-strong prominence relation is preserved in both cases.

Perceptually, prominence as a redundant system facilitates robust communication by encompassing multiple cues serving the same function. Listeners have been shown to be able to make use of a variety of cues to decode prominence in many previous studies \citep{turnbullProminencePerceptionDependent2017, baumannWhatMakesWord2018, coleSoundStructureMeaning2019, wagnerExploitingSpeechgestureLink2019}. Our own perception study presented in Chapter \ref{ch: Perception} confirms these findings. This mechanism ensures communicative success despite the general lack of invariance in prosody-meaning mapping \citep{winterSpokenLanguageAchieves2014}. As argued in \citet{baumannWhatMakesWord2018}, listeners need not pay attention to every single cue but can compensate for their inattentiveness to certain parts of the signal by relying on other cues. In listeners who are physically able to perceive and make use of all available acoustic dimensions, such compensatory effects may simply reflect personal preferences for certain cues over others. However, for individuals with perceptual deficits such as congenital amusia, which severely limits pitch perception, functional redundancy crucially enables them to develop successful speech perception strategies by compensating for a missing cue \citep{jasminTailoredPerceptionIndividuals2020}.

Nevertheless, some prominence cues are generally preferred to others. Especially the categorical prosodic variables accent status and accent type are strong cues to prominence. Prosodic categories can be regarded as inherently redundant, since they encompass several acoustic cues (see also previous section). Accent status also conveys the relationality of prominence, as it is the most direct reflection of metrical structure. That is, the strength of the prenuclear or nuclear accent derives from its place in the hierarchical organization of prosodic constituents \citep{laddIntonationalPhonology2008, calhounCentralityMetricalStructure2010}. Naturally, this makes accent status the strongest cue to prominence not only on a general but also at the individual level.

From a production perspective, redundancy enables inter-individual variability. There are many different ways to encode prominence relations and individuals can and do exploit this functional redundancy. The production studies presented in Chapters \ref{ch:Production I} and \ref{ch:Production II} provide evidence that speakers rarely make use of all prosodic parameters available to them to encode a meaning contrast. Instead, they choose different subsets of cues to mark different functional aspects in a compensatory manner. That is, although there is usually some redundancy in the productions of individual speakers, most speakers are not maximally redundant themselves. Redundancy thus at the same time ensures robustness in perception and enables variability in production.

\section{Limitations of the present work and directions for future research}

There are some limitations to the studies presented in this book. Firstly, all production data represent read speech and the perception data is acquired using excerpts of read speech as stimuli. Some differences between read speech and spontaneous speech styles are known. Characteristics of read speech include higher median F0 \citep[English:][]{hazanDoesReadingClearly2010}, steeper declination, and stronger F0 resetting \citep[Swedish:][]{swertsF0DeclinationReadaloud1996}. \citet{baumannCoreferenceLexicalGivenness2013} for German found more early peaks in read speech and more low nuclear accents followed by high boundaries (i.e., continuation rises) in spontaneous speech. Generally, spontaneous and read speech thus appear to differ systematically. However, \citet{grylliaManyShapes2022} observed some variability with regard to the extent that Greek speakers differed between speaking styles: Some speakers produced more animated speech in a reading task, while others did not vary between read versus unscripted speech. Therefore, the distinction between lively and monotonous speech is potentially more relevant for prosody than scriptedness. Lively speech is characterized by more nuclear and more rising accents, steeper F0 excursions. and larger ranges \citep{baumannInformativenessSpeakingStyle2019, chodroffInformationStructureAffect2018}. In our production studies, we elicited such lively speech via listener-directed tasks. Nevertheless, there may be other disadvantages to read speech: \citet{bardDissociationDeaccentingGivenness1999} observed that deaccentuation of second mentions and repetitions rarely occurred within sentences of the same structure in unscripted speech. Spontaneous data could thus offer new insights into the encoding of information status that are not available in scripted speech. 

The presented studies are exclusively concerned with the spoken modality. However, gestures may additionally function as visual cues to prominence \citep{krahmerEffectsVisualBeats2007, ambrazaitisMultimodalProminencesExploring2017, esteve-gibertTimingHeadMovements2017, rohrerVisualizingProsodicStructure2023}. As such, visual cues add a further dimension of redundancy to prominence, which can be exploited by individuals.
Visual prominence cues may potentially be used as a compensatory strategy by individuals not attending to acoustic cues of prominence. Initial evidence from perception showing that individuals rely on visual cues to different extents comes from \citet{bujokAudiovisualPerceptionLexical2024}. Furthermore, it has been shown that individuals vary with regard to how well they integrate auditory and visual cues \citep[e.g.,][]{wilbiksSpeechNonspeechMeasures2022}. However, production studies explicitly  concerned with individual differences in the use of gesture are currently lacking.

Another interesting remaining question is that of a potential link between production and perception. A relation between robust prosodic encoding and an enhanced ability to decode prosodically encoded contrasts in the same individual could help provide a holistic account of individual variability in prosody. Although we can certainly expect intra-individual variability with respect to which contrasts speakers encode most robustly and decode most proficiently, the production-perception link should remain constant and reflect an individuals' general aptitude for both the encoding and the decoding of a certain meaning contrast. Since different participants were included in the production and perception studies, the present book cannot contribute to this line of research. However, evidence for a production-perception link in previous studies is inconclusive so far (see Section \ref{sec:ind_perception}), so that future studies in this direction need to take an exploratory approach.

Throughout all studies, we have noted sociodemographic characteristics associated with speakers and listeners employing different strategies. No clear patterns emerged with regard to age, gender, and regional effects on prominence production and perception. We observed no systematic age effects, which can possibly be attributed to relatively constrained age ranges in our participants, particularly in the production studies. With regard to gender, the strong encoders of lexical givenness in Production Study II were predominantly female. However, this pattern was not confirmed in the encoding of the other contrasts in the same study nor in the other two studies. Similarly, regional distributions did not systematically map onto differences in prominence encoding and decoding. Therefore, we assume that any individual variability observed in the present data is independent of the sociodemographic characteristics of participants and truly idiosyncratic. However, studies with more controlled, properly stratified participant samples could yield further insights into this issue. 

In general, the observations reported in the present book need to be extended beyond laboratory studies to more natural speech in production and communicatively relevant tasks in perception. More natural speech almost necessarily entails speech in face-to-face interaction accompanied by the visual modality, which offers a further avenue of investigation into the redundant nature of prominence. In addition, more differentiated measures of personality traits and abilities could offer new insights into underlying reasons for individual variability, which is likely to also emerge in such settings.

\section{Conclusion} \label{sec: conclusions}

The present book contributes to the study of inter-individual differences in prominence production and perception, which has recently seen considerable interest in the speech prosody community \textendash{} and for good reasons. Individual variability emerged in all three studies presented here, although speakers and listeners were not necessarily drastically different from a group-level average. This is reassuring for many studies not explicitly considering the individual. Nevertheless, acknowledging and describing individual differences is necessary to achieve empirical accuracy since prosodic form-meaning mapping rarely looks the same for or is interpreted in exactly the same way by every single individual. This caveat is especially relevant in light of the drive for ever larger data sets, perpetuating the risk of losing sight of the individual. Understanding individual variability also allows us to refine theoretical assumptions and provides new perspectives on already established concepts in prosodic theory, such as the dichotomy of tune and metrical structure, relationality and redundancy in prominence encoding and decoding, and probabilistic form-meaning mapping. The present work has shown potential avenues into this large endeavor. 